\clearpage\begingroup\SourceLeaf{82b}\HistoricalHead{77}{R}{GAUSS AN GERLING}
\begin{center}[4.]\end{center}
\begin{center}\textsc{Gauss} an \textsc{Gerling}.\quad Göttingen, 2. April 1832.\end{center}
\noindent . . . . . . . . . . . . . . . . . . . . .\par
\indent Ich habe seit etwa einem Monat mich recht viel mit dem Magnetismus
beschäftigt, und angefangen, nicht bloss diejenigen Ideen, die ich Ihnen Weihnachten
mittheilte, selbst (unter vielfachem Beistande von Freund \textsc{Weber}) auszuführen,
sondern alles noch viel weiter auszudehnen. Ich komme fast täglich
noch auf eine neue Idee, und muss nur bedauern, dass die Ausführung, wobei
anfangs bald dies bald jenes erst weitläufig herbeigeschafft werden muss, nicht
so schnell damit Schritt halten kann. Aber auch wie es jetzt steht, ist meine
Erwartung weit übertroffen. Die tägliche Variation kann ich schon fast von
Minute zu Minute verfolgen und wenige Bogensekunden (sage z. B. 2 oder 3)
\textls[120]{sicher} sichtbar machen. Ich hoffe in allen einzelnen Momenten, nemlich
Intensität, Declination, Inclination und Variation dieser drei Elemente die bisherige
Schärfe weit überbieten zu können. Die Schwingungsdauer bestimme
ich schon jetzt mit fast unglaublicher Schärfe. . . . . . . . . . . . . . . . . . .
\clearpage\endgroup
